\documentclass[12pt, a4paper]{article}
\usepackage[UTF8]{ctex}
\usepackage{amsmath, amssymb, amsthm}
\usepackage{geometry}
\usepackage{graphicx}
\usepackage{bm}

\geometry{left=2.5cm, right=2.5cm, top=2.5cm, bottom=2.5cm}

\title{\textbf{计算机视觉 作业3：三维重建} \\ 第一部分：理论}
\author{许政阳 \\ 学号：231880081}
\date{}

\begin{document}
\maketitle

\section*{问题 1：证明题}

\subsection*{Q1.1 [10 分]}

\textbf{题目：}假设两个相机都注视空间中的同一个三维点 $\mathbf{X}$，即它们的主轴相交于该点。若该三维点在两张图像中的投影都位于像素坐标 $(0, 0)$，证明基础矩阵 $F$ 的 $f_{33}$ 元素必须为 $0$。

\bigskip
\textbf{证明：}

基础矩阵 $F$ 定义了对极约束：
\[
\mathbf{x}'^{\top} F \, \mathbf{x} = 0
\]
其中 $\mathbf{x}$ 和 $\mathbf{x}'$ 分别为三维点 $\mathbf{X}$ 在左、右图像中的齐次坐标。

由题意，三维点 $\mathbf{X}$ 在两张图像中的投影都位于像素坐标 $(0,0)$，因此其齐次坐标为：
\[
\mathbf{x} = \begin{pmatrix} 0 \\ 0 \\ 1 \end{pmatrix}, \quad
\mathbf{x}' = \begin{pmatrix} 0 \\ 0 \\ 1 \end{pmatrix}
\]

将基础矩阵 $F$ 展开为各元素：
\[
F = \begin{pmatrix}
f_{11} & f_{12} & f_{13} \\
f_{21} & f_{22} & f_{23} \\
f_{31} & f_{32} & f_{33}
\end{pmatrix}
\]

代入对极约束：
\[
\mathbf{x}'^{\top} F \, \mathbf{x} =
\begin{pmatrix} 0 & 0 & 1 \end{pmatrix}
\begin{pmatrix}
f_{11} & f_{12} & f_{13} \\
f_{21} & f_{22} & f_{23} \\
f_{31} & f_{32} & f_{33}
\end{pmatrix}
\begin{pmatrix} 0 \\ 0 \\ 1 \end{pmatrix}
\]

首先计算 $F \, \mathbf{x}$：
\[
F \, \mathbf{x} =
\begin{pmatrix}
f_{11} \cdot 0 + f_{12} \cdot 0 + f_{13} \cdot 1 \\
f_{21} \cdot 0 + f_{22} \cdot 0 + f_{23} \cdot 1 \\
f_{31} \cdot 0 + f_{32} \cdot 0 + f_{33} \cdot 1
\end{pmatrix}
= \begin{pmatrix} f_{13} \\ f_{23} \\ f_{33} \end{pmatrix}
\]

再计算 $\mathbf{x}'^{\top} (F \, \mathbf{x})$：
\[
\mathbf{x}'^{\top} F \, \mathbf{x} =
\begin{pmatrix} 0 & 0 & 1 \end{pmatrix}
\begin{pmatrix} f_{13} \\ f_{23} \\ f_{33} \end{pmatrix}
= 0 \cdot f_{13} + 0 \cdot f_{23} + 1 \cdot f_{33} = f_{33}
\]

由于对极约束要求 $\mathbf{x}'^{\top} F \, \mathbf{x} = 0$，因此：
\[
\boxed{f_{33} = 0}
\]

\hfill $\blacksquare$

\newpage

\subsection*{Q1.2 [10 分]}

\textbf{题目：}假设一个机器人携带惯性传感器运动，传感器在时刻 $i$ 报告机器人在全局坐标系中的绝对旋转 $R_i$ 和平移 $t_i$。求时刻 $i$ 到 $i+1$ 之间的相对旋转 $R_{\text{rel}}$ 和相对平移 $t_{\text{rel}}$。若已知相机内参 $K$，请用 $K$、$R_{\text{rel}}$ 和 $t_{\text{rel}}$ 表示本质矩阵 $E$ 与基础矩阵 $F$。

\bigskip
\textbf{解答：}

\subsubsection*{第一部分：求 $R_{\text{rel}}$ 和 $t_{\text{rel}}$}

传感器在时刻 $i$ 报告的全局位姿 $(R_i, t_i)$ 表示从相机坐标系到全局坐标系的变换。对于全局坐标系中的一个三维点 $\mathbf{P}_w$，其在相机 $i$ 坐标系中的表示为：
\[
\mathbf{P}_i = R_i^{\top} (\mathbf{P}_w - t_i) = R_i^{\top} \mathbf{P}_w - R_i^{\top} t_i
\]

类似地，在相机 $i+1$ 坐标系中：
\[
\mathbf{P}_{i+1} = R_{i+1}^{\top} (\mathbf{P}_w - t_{i+1}) = R_{i+1}^{\top} \mathbf{P}_w - R_{i+1}^{\top} t_{i+1}
\]

为了建立从相机 $i$ 到相机 $i+1$ 的相对变换，我们需要将 $\mathbf{P}_{i+1}$ 用 $\mathbf{P}_i$ 表示。

从相机 $i$ 坐标系恢复全局坐标：
\[
\mathbf{P}_w = R_i \, \mathbf{P}_i + t_i
\]

代入相机 $i+1$ 的表达式：
\begin{align}
\mathbf{P}_{i+1} &= R_{i+1}^{\top} (R_i \, \mathbf{P}_i + t_i - t_{i+1}) \nonumber \\
&= R_{i+1}^{\top} R_i \, \mathbf{P}_i + R_{i+1}^{\top} (t_i - t_{i+1})
\end{align}

因此，相对旋转和相对平移为：
\[
\boxed{R_{\text{rel}} = R_{i+1}^{\top} R_i}
\]
\[
\boxed{t_{\text{rel}} = R_{i+1}^{\top} (t_i - t_{i+1})}
\]

\subsubsection*{第二部分：用 $K$、$R_{\text{rel}}$、$t_{\text{rel}}$ 表示 $E$ 和 $F$}

\textbf{本质矩阵 $E$：}

本质矩阵编码了两个相机之间的对极几何关系。根据 Longuet-Higgins 方程，对于归一化相机坐标下的对应点 $\hat{\mathbf{x}}$ 和 $\hat{\mathbf{x}}'$，有：
\[
\hat{\mathbf{x}}'^{\top} E \, \hat{\mathbf{x}} = 0
\]

本质矩阵由相对旋转和相对平移定义为：
\[
\boxed{E = [t_{\text{rel}}]_{\times} \, R_{\text{rel}}}
\]
其中 $[t_{\text{rel}}]_{\times}$ 是 $t_{\text{rel}}$ 的反对称矩阵（skew-symmetric matrix）：若 $t_{\text{rel}} = (t_x, t_y, t_z)^{\top}$，则
\[
[t_{\text{rel}}]_{\times} = \begin{pmatrix}
0 & -t_z & t_y \\
t_z & 0 & -t_x \\
-t_y & t_x & 0
\end{pmatrix}
\]

\textbf{基础矩阵 $F$：}

基础矩阵是本质矩阵在像素坐标系下的推广。图像点 $\mathbf{x}$ 与归一化相机坐标 $\hat{\mathbf{x}}$ 之间的关系为：
\[
\mathbf{x} = K \, \hat{\mathbf{x}} \quad \Longrightarrow \quad \hat{\mathbf{x}} = K^{-1} \mathbf{x}
\]

将 $\hat{\mathbf{x}} = K^{-1} \mathbf{x}$ 和 $\hat{\mathbf{x}}' = K^{-1} \mathbf{x}'$ 代入本质矩阵约束：
\begin{align}
(K^{-1} \mathbf{x}')^{\top} E \, (K^{-1} \mathbf{x}) &= 0 \nonumber \\
\mathbf{x}'^{\top} K^{-\top} E \, K^{-1} \mathbf{x} &= 0 \nonumber \\
\mathbf{x}'^{\top} F \, \mathbf{x} &= 0
\end{align}

因此，基础矩阵为：
\[
\boxed{F = K^{-\top} E \, K^{-1} = K^{-\top} [t_{\text{rel}}]_{\times} \, R_{\text{rel}} \, K^{-1}}
\]

\hfill $\blacksquare$

\end{document}
